% !PRINTED 203
\weberpage{203}

%<<CONT
besteht, so kann darin die Variable $v$ durch $v + 4\mu K + 4\mu' iK'$
ersetzt werden. Wegen der Periodizität von sn\,$v$, cn\,$v$, dn\,$v$ folgt
aber daraus, daß die Gleichung

\[
F(x,\ \mathrm{sn}\,v,\ \mathrm{cn}\,v,\ \mathrm{dn}\,v) = 0
\]

für \gesperrt{alle} Wurzeln der Gleichung (4) erfüllt ist. Um ihre \gesperrt{Galois\-
sche} Gruppe zu ermitteln, setzen wir zunächst den \gesperrt{Rationali\-
tätsbereich} fest. Er soll folgende Größen umfassen:

\begin{enumerate}
\item[1.] rationale Zahlen,
\item[2.] rationale Funktionen von $\varkappa^2$,
\item[3.] die drei Funktionen sn\,$v$, cn\,$v$, dn\,$v$,
\item[4.] die Größen $\mathrm{sn}\left(\dfrac{4\mu K' + 4\mu' iK'}{n}\right)$. % [sic]
\end{enumerate}

Aus (2) folgt dann, daß auch

\[
\mathrm{cn}\left(\frac{4\mu K + 4\mu' iK'}{n}\right), \qquad \mathrm{dn}\left(\frac{4\mu K + 4\mu' iK'}{n}\right)
\]

zum Rationalitätsbereich gehören.

Wir bemerken nun, daß nach dem Additionstheorem jede der
Größen $x_{\mu,\mu'}$ durch jede andere unter ihnen rational ausdrückbar
ist. Wenn insbesondere

\weq{5}{x_{\mu,\mu'} = F_{\mu,\mu'}(x_{0,0})}

ist, so ist

\weq{6}{x_{\mu+\nu,\mu'+\nu'} = F_{\mu,\mu'}(x_{\nu,\nu'}) = F_{\mu+\nu,\mu'+\nu'}(x_{0,0}).}

Daraus folgt, daß die \gesperrt{Galoissche} Gruppe unserer Gleichung
aus den $n^2$ Vertauschungen $S_{\nu,\nu'}$ besteht, die man erhält, wenn
man in den Wurzeln $x_{\mu,\mu'}$ die Indizes $\mu, \mu'$ alle um dasselbe
Zahlenpaar $\nu, \nu'$ vermehrt (wobei jeder Index nach dem Modul $n$
zu nehmen ist).

Denn nach (5) kann jede rationale Funktion der Wurzeln
$x_{\mu,\mu'}$ rational durch $x_{0,0}$ in der Form

\[
\Phi(x_{0,0})
\]

ausgedrückt werden und geht also nach (6) durch die Substitution
$S_{\nu,\nu'}$ in $\Phi(x_{\nu,\nu'})$ über; bleibt sie also durch diese Substitution
ungeändert, so ist \gesperrt{sie rational}. Umgekehrt folgt aus der
Irreduzibilität, daß jede rationale Gleichung zwischen den Wurzeln,
da sie in die Form gesetzt werden kann:

\[
\psi(x_{0,0}) = 0,
\]

und also

\[
\psi(x_{\nu,\nu'}) = 0
\]
%CONT>>
